% ID: FIS-FLU-PID-004
% Titolo: Getti d'acqua da fori a profondita diverse
% Disciplina: Fisica
% Area: Fluidostatica
% Argomento: Pressione idrostatica
% Sottoargomento: Effetto della profondita sulla velocita di uscita
% Classe: Terza liceo artistico
% Difficolta: 3
% Tipo: quesito_teorico
% Risultato: soluzione da aggiungere
% Tempo_stimato: 8 min
% Competenze: interpretazione fisica, collegamento tra pressione e moto, argomentazione
% Prerequisiti: pressione idrostatica, profondita, moto parabolico
% Tag: fisica, fluidostatica, pressione_idrostatica, getti_acqua, profondita
% Fonte: verifica_fisica_fluidostatica_anonimizzata
% Autore: Riccardo Carli
% Licenza: CC-BY-SA-4.0

\begin{esercizio}
Un recipiente verticale pieno d'acqua ha tre piccoli fori sulla stessa parete,
posti a profondita diverse rispetto alla superficie libera. Quando i fori
vengono aperti, dai tre fori escono getti d'acqua con traiettorie diverse.

Spiega perche i getti non hanno la stessa forma e formula un'ipotesi sul tipo
di curva seguita da ciascun getto.
\end{esercizio}

\begin{soluzione}
% Soluzione da aggiungere.
\end{soluzione}

